% =============================================================================
% Arena Limpia - Reporte del proyecto final (español)
% Laboratorio de Computación Gráfica e Interacción Humano-Computadora, FI UNAM, 2027-1
%
% Documento vivo: se va llenando conforme avanza el proyecto. Lo marcado con \pendiente{}
% falta por completar antes de la entrega.
%
% Compilar desde esta carpeta (dos veces, para referencias cruzadas):
%     pdflatex reporte.tex && pdflatex reporte.tex
% =============================================================================
\documentclass[12pt,letterpaper]{article}

\usepackage[utf8]{inputenc}
\usepackage[T1]{fontenc}
\usepackage[spanish,es-tabla,es-nodecimaldot]{babel}
\usepackage{lmodern}
\usepackage[margin=2.5cm]{geometry}
\usepackage{setspace}
\onehalfspacing                         % interlineado 1.5 (requisito)
\usepackage{graphicx}
\usepackage{float}
\usepackage{caption}
\usepackage{subcaption}
\usepackage{booktabs}
\usepackage{tabularx}
\usepackage{amsmath}
\usepackage{xcolor}
\usepackage{listings}
\usepackage{tikz}
\usetikzlibrary{arrows.meta,positioning}
\usepackage[hidelinks]{hyperref}

\graphicspath{{figuras/}}

% Texto pendiente: se ve en rojo para no olvidarlo
\newcommand{\pendiente}[1]{\textcolor{red}{[Pendiente: #1]}}

\lstset{
  basicstyle=\ttfamily\footnotesize,
  keywordstyle=\color{blue!60!black}\bfseries,
  commentstyle=\color{green!40!black}\itshape,
  stringstyle=\color{orange!70!black},
  numbers=left, numberstyle=\tiny\color{gray},
  frame=single, framerule=0.3pt, rulecolor=\color{gray!60},
  breaklines=true, showstringspaces=false, tabsize=4,
  inputencoding=utf8, extendedchars=true,
  literate={á}{{\'a}}1 {é}{{\'e}}1 {í}{{\'i}}1 {ó}{{\'o}}1 {ú}{{\'u}}1
           {ñ}{{\~n}}1 {Á}{{\'A}}1 {É}{{\'E}}1 {Ó}{{\'O}}1 {ü}{{\"u}}1
}
\renewcommand{\lstlistingname}{Código}
\renewcommand{\lstlistlistingname}{Índice de códigos}

\begin{document}

% -----------------------------------------------------------------------------
% Portada
% -----------------------------------------------------------------------------
\begin{titlepage}
  \centering
  {\large Universidad Nacional Autónoma de México\par}
  {\large Facultad de Ingeniería\par}
  \vspace{0.5cm}
  {\normalsize Laboratorio de Computación Gráfica e Interacción Humano-Computadora\par}
  {\normalsize Semestre 2027-1\par}
  \vspace{2.5cm}
  {\Huge\bfseries Arena Limpia\par}
  \vspace{0.3cm}
  {\Large Torneo de Voleibol de Playa\par}
  \vspace{0.5cm}
  {\large Proyecto final\par}
  \vfill
  \begin{tabular}{l l}
    \textbf{Autores:} & Hernández Yukioayax Gabriel \\
                      & Parra Bello Carlos Enrique \\[0.3cm]
    \textbf{Contacto:} & \pendiente{correos institucionales} \\[0.3cm]
    \textbf{Profesor:} & \pendiente{nombre completo del profesor} \\[0.3cm]
    \textbf{Fecha de entrega:} & \pendiente{fecha}
  \end{tabular}
  \vspace{1.5cm}
\end{titlepage}

% -----------------------------------------------------------------------------
% Hoja de evidencias
% -----------------------------------------------------------------------------
\section*{Hoja de evidencias}
\begin{tabular}{l l}
  \textbf{Título del trabajo:} & Arena Limpia: Torneo de Voleibol de Playa \\[0.3cm]
  \textbf{Participantes} & \textbf{Número de cuenta} \\
  Hernández Yukioayax Gabriel & \pendiente{número de cuenta} \\
  Parra Bello Carlos Enrique  & \pendiente{número de cuenta} \\[0.3cm]
  \textbf{Video demostrativo:} & \pendiente{enlace al video}
\end{tabular}

\newpage
\tableofcontents
\newpage

% -----------------------------------------------------------------------------
\section*{Resumen}
\addcontentsline{toc}{section}{Resumen}
% -----------------------------------------------------------------------------
Arena Limpia es un sistema gráfico interactivo desarrollado en C++ con OpenGL 3.3 y GLSL que
representa un torneo de voleibol de playa en una costa. La escena reúne modelos geométricos
propios y de otros autores colocados a escala real, un mapa de ambiente tipo \emph{cube map}
obtenido de una fotografía panorámica de playa, iluminación de Phong con una fuente
direccional, dos modos de cámara (libre y en tercera persona) y tres tipos de animación:
animaciones con esqueleto basadas en cuadros clave (caminar, saque, remate, festejo y la señal
del árbitro), y animaciones procedurales como el tiro parabólico del balón con rebotes y
choque con la red. El usuario controla a un jugador, camina por la arena y ejecuta las
jugadas con el teclado.

\noindent\textbf{Palabras clave:} computación gráfica, OpenGL, animación esquelética,
\emph{skinning}, cuadros clave, animación procedural, tiro parabólico, voleibol de playa.

% -----------------------------------------------------------------------------
\section{Introducción}
% -----------------------------------------------------------------------------
La computación gráfica interactiva combina modelado geométrico, iluminación, texturizado y
animación para producir imágenes en tiempo real a partir de datos tridimensionales. El
modelo de iluminación de Phong \cite{phong1975} sigue siendo la base de muchos sistemas en
tiempo real por su bajo costo y su buena aproximación a materiales opacos y plásticos. En
animación de personajes, la técnica dominante es la animación esquelética con
\emph{linear blend skinning} \cite{magnenat1988}: cada vértice de la malla se deforma con un
promedio ponderado de las transformaciones de los huesos que lo afectan, lo que permite
animar mallas complejas moviendo pocos huesos. Las poses de los huesos se definen en cuadros
clave e interpolando entre ellos, una práctica heredada de la animación tradicional
\cite{lasseter1987}. Por último, la animación procedural genera el movimiento con reglas o
leyes físicas en lugar de poses guardadas; el ejemplo clásico es el tiro parabólico, que
describe el vuelo de un objeto bajo gravedad constante \cite{halliday2014}.

El proyecto aplica estas técnicas a un tema deportivo: un torneo de voleibol de playa con su
cancha, gradas, mesa de jueces, premiación y la vida de la playa alrededor (puesto de
helados, salvavidas, fauna y embarcaciones). El nombre \emph{Arena Limpia} hace referencia al
cuidado de la playa, presente en la escena con botes de basura y objetos reciclables.

\pendiente{estado del arte: comparar con otros simuladores o videojuegos de voleibol de playa
y con escenas de playa en motores gráficos, con citas.}

\pendiente{análisis económico, social y ético del tema.}

% -----------------------------------------------------------------------------
\section{Metodología}
% -----------------------------------------------------------------------------
El desarrollo siguió las etapas de planteamiento del problema, análisis, diseño, codificación
y pruebas. El código vive en un repositorio Git compartido por el equipo; cada cambio se
probó compilando el programa y generando capturas de pantalla automáticas desde la línea de
comandos (sección~\ref{sec:pruebas}).

\subsection{Requerimientos}
\begin{table}[H]
  \centering
  \caption{Requerimientos de hardware y software.}
  \label{tab:requerimientos}
  \begin{tabularx}{\linewidth}{l X}
    \toprule
    \textbf{Elemento} & \textbf{Requerimiento} \\
    \midrule
    Tarjeta gráfica & Compatible con OpenGL 3.3 o superior \\
    Sistema operativo & Windows 10 u 11 de 64 bits \\
    Memoria & \pendiente{medir uso de memoria} \\
    Desarrollo & Visual Studio 2026 (C++17), GLFW, GLAD, GLM, Assimp, stb\_image \\
    Modelado y animación & Blender 5.2 con scripts de Python \\
    \bottomrule
  \end{tabularx}
\end{table}

\subsection{Arquitectura del programa}
El programa es un ciclo de juego: en cada cuadro lee la entrada del usuario, avanza la
simulación (balón, personajes y árbitro) y dibuja la escena (figura~\ref{fig:arquitectura}).

\begin{figure}[H]
  \centering
  \begin{tikzpicture}[
      node distance=0.6cm and 0.9cm,
      caja/.style={draw, rounded corners, align=center, minimum height=1cm,
                   text width=3.1cm, font=\small},
      flecha/.style={-{Stealth[length=2.5mm]}, thick}]
    \node[caja, fill=blue!8] (inicio) {Inicio: carga de shaders, skybox y modelos};
    \node[caja, fill=orange!10, right=of inicio] (entrada) {Entrada: teclado y mouse};
    \node[caja, fill=green!10, right=of entrada] (sim) {Simulación: balón, personajes, árbitro};
    \node[caja, fill=purple!10, below=of sim] (render) {Dibujo: skybox, arena, modelos};
    \node[caja, fill=gray!10, left=of render] (swap) {Intercambio de búferes};
    \draw[flecha] (inicio) -- (entrada);
    \draw[flecha] (entrada) -- (sim);
    \draw[flecha] (sim) -- (render);
    \draw[flecha] (render) -- (swap);
    \draw[flecha] (swap) -- (entrada);
  \end{tikzpicture}
  \caption{Ciclo principal del programa (\texttt{src/main.cpp}).}
  \label{fig:arquitectura}
\end{figure}

\begin{table}[H]
  \centering
  \caption{Archivos principales del código.}
  \label{tab:archivos}
  \begin{tabularx}{\linewidth}{l X}
    \toprule
    \textbf{Archivo} & \textbf{Contenido} \\
    \midrule
    \texttt{src/main.cpp} & Escena, cámaras, entrada, simulación del balón y ciclo principal \\
    \texttt{include/skinnedmodel.h} & Carga de modelos con esqueleto y cálculo de poses \\
    \texttt{include/model.h}, \texttt{mesh.h} & Carga y dibujo de modelos OBJ (base del curso) \\
    \texttt{shaders/phong.vs}, \texttt{phong.fs} & Iluminación de Phong con el sol \\
    \texttt{shaders/phong\_skinning.vs} & Deformación de la malla con los huesos \\
    \texttt{tools/animar\_personajes.py} & Creación y exportación de las animaciones en Blender \\
    \texttt{tools/exportar\_obj.py} & Exportación de modelos de Blender a OBJ \\
    \bottomrule
  \end{tabularx}
\end{table}

\subsection{Modelado y escala}
\label{sec:escala}
Los modelos se hicieron en Blender y se exportan a OBJ. Cada integrante modeló con sus
propias unidades: un personaje medía 10.7 unidades de alto, una botella 0.09 y un bote 558,
así que no se podía usar una escala única. La solución fue describir cada objeto de la escena
por su \textbf{altura real en metros}; al cargarlo se calcula su caja envolvente y la escala
es
\begin{equation}
  s = \frac{h_{\text{real}}}{y_{\max} - y_{\min}},
\end{equation}
y el objeto se traslada para que la base de su caja quede sobre la arena. Así, por ejemplo,
los jugadores miden entre 1.72 y 1.85~m, el balón 0.21~m y el faro 20~m. La red se escala de
forma no uniforme para que sus postes coincidan con el ancho de la cancha.

\subsection{Cámaras}
Se implementaron dos modelos de cámara; la tecla \textbf{V} cambia entre ellos.
\begin{itemize}
  \item \textbf{Cámara libre:} se mueve con W, A, S, D, sube y baja con Espacio y C, y mira con
        el mouse (ángulos de guiñada y cabeceo).
  \item \textbf{Tercera persona:} sigue a uno de los cuatro jugadores (teclas 1 a 4). La cámara
        orbita alrededor del jugador a una distancia ajustable con la rueda del mouse:
        \begin{equation}
          \mathbf{e} = \mathbf{t} - d\,(\cos\phi\cos\theta,\ \sin\phi,\ \cos\phi\sin\theta),
        \end{equation}
        donde $\mathbf{t}$ es un punto a la altura de los hombros del jugador, $d$ la
        distancia, $\theta$ la guiñada y $\phi$ el cabeceo. W, A, S, D mueven al jugador en la
        dirección en la que mira la cámara y el jugador gira gradualmente hacia donde camina.
\end{itemize}

\subsection{Animación esquelética}
\label{sec:esqueleto}
Los once personajes comparten un esqueleto de 19 huesos (cadera, columna, pecho, cuello,
cabeza, y hombro, brazo, antebrazo, mano, muslo, pierna y pie de cada lado). El formato OBJ
no guarda huesos, así que los personajes animados se exportan a \textbf{glTF 2.0}
\cite{gltf2021}, que guarda la malla, los pesos de cada vértice, el esqueleto y las
animaciones. Assimp \cite{assimp} lo carga y la clase \texttt{SkinnedModel} calcula en cada
cuadro la matriz de cada hueso:
\begin{equation}
  M_i = G^{-1}\, N_i(t)\, O_i,
\end{equation}
donde $N_i(t)$ es la transformación global del hueso en el instante $t$ (producto de las
transformaciones de sus padres interpoladas entre cuadros clave), $O_i$ la matriz que lleva
la malla al espacio del hueso en reposo y $G^{-1}$ la inversa de la transformación de la
raíz. El \emph{vertex shader} deforma cada vértice con hasta cuatro huesos
(código~\ref{lst:skinning}):
\begin{equation}
  \mathbf{v}' = \sum_{k=1}^{4} w_k\, M_{j_k}\, \mathbf{v}, \qquad \sum_k w_k = 1.
\end{equation}

\begin{lstlisting}[language=C, caption={Deformación de vértices con los huesos (\texttt{shaders/phong\_skinning.vs}).}, label={lst:skinning}]
mat4 skin = gBones[int(bIDs1.x)] * bWeights1.x
          + gBones[int(bIDs1.y)] * bWeights1.y
          + gBones[int(bIDs1.z)] * bWeights1.z
          + gBones[int(bIDs1.w)] * bWeights1.w;
vec4 worldPos = model * skin * vec4(aPos, 1.0);
Normal = mat3(transpose(inverse(model * skin))) * aNormal;
\end{lstlisting}

La clase \texttt{SkinnedModel} se basó en la clase \texttt{AnimatedModel} de la práctica de
animación del curso, con cambios necesarios para este proyecto: numerar los huesos por
nombre para todo el modelo, calcular normales cuando el archivo no las trae, conservar la
transformación de los nodos sin animación, elegir la animación por nombre y medir el tiempo
en segundos.

\subsection{Animaciones por cuadros clave}
Las animaciones se crean con el script \texttt{tools/animar\_personajes.py}, que abre el
archivo de Blender de cada personaje sin modificarlo, define las poses en ciertos cuadros y
exporta el resultado. Las rotaciones se escriben en los ejes del esqueleto y se convierten al
eje local de cada hueso,
\begin{equation}
  q_{\text{local}} = r^{-1}\, q_{\text{esqueleto}}\, r,
\end{equation}
donde $r$ es la orientación del hueso en reposo. En cada cuadro clave la cadera se ajusta
para que el pie más bajo toque el piso, y los saltos se suman en metros; así los personajes
no flotan ni se hunden. La tabla~\ref{tab:keyframes} resume las animaciones.

\begin{table}[H]
  \centering
  \caption{Animaciones por cuadros clave (24 cuadros por segundo).}
  \label{tab:keyframes}
  \begin{tabularx}{\linewidth}{l l l X}
    \toprule
    \textbf{Animación} & \textbf{Duración} & \textbf{Personajes} & \textbf{Cómo se activa} \\
    \midrule
    Caminar  & 1.00 s (ciclo) & jugadores, árbitro & Moverse en tercera persona; su velocidad depende de la del jugador \\
    Saque    & 1.67 s & jugadores & Tecla Q \\
    Remate   & 1.75 s & jugadores & Tecla E (incluye un salto de 0.6 m) \\
    Festejo  & 2.42 s & jugadores & Tecla F \\
    Marcar   & 1.58 s & árbitro   & Automática, cada 8 s \\
    Parado   & 1.00 s (ciclo) & todos los animados & En reposo \\
    \bottomrule
  \end{tabularx}
\end{table}

Para que los pies no se resbalen, la animación de caminar avanza al ritmo de los pasos: un
ciclo cubre 1.5~m, así que su velocidad de reproducción es $v / 1.5$ ciclos por segundo.

\subsection{Animaciones procedurales}
\subsubsection{Tiro parabólico del balón}
\label{sec:tiro}
En el saque el balón sigue la mano izquierda del jugador; su posición se obtiene del hueso
\texttt{mano.L} en la pose actual. Al lanzarlo, y otra vez al golpearlo, se calcula la
velocidad inicial $\mathbf{v}_0$ que lleva el balón de $\mathbf{p}_0$ a $\mathbf{p}_1$ en un
tiempo $T$, despejando la ecuación del tiro parabólico:
\begin{equation}
  \mathbf{p}_1 = \mathbf{p}_0 + \mathbf{v}_0 T + \tfrac{1}{2}\mathbf{g}T^2
  \quad\Longrightarrow\quad
  \mathbf{v}_0 = \frac{\mathbf{p}_1 - \mathbf{p}_0}{T} - \tfrac{1}{2}\mathbf{g}T.
\end{equation}
En el lanzamiento, $\mathbf{p}_1$ es la posición que tendrá la mano derecha en el cuadro del
golpe, así que el balón llega a la mano exactamente cuando la animación lo golpea. En el
remate el balón llega como un pase desde un costado. Después del golpe, el balón se integra
cada cuadro con la solución exacta para aceleración constante,
\begin{equation}
  \mathbf{p} \leftarrow \mathbf{p} + \mathbf{v}\,\Delta t + \tfrac{1}{2}\mathbf{g}\,\Delta t^2,
  \qquad \mathbf{v} \leftarrow \mathbf{v} + \mathbf{g}\,\Delta t,
\end{equation}
y se resuelven tres contactos: con la arena rebota conservando el 55\% de la velocidad
vertical y el 75\% de la horizontal, y rueda frenando hasta detenerse; si cruza el plano de la
red a la altura de la red, rebota hacia atrás. El balón gira alrededor del eje perpendicular a
su velocidad horizontal con velocidad angular $\omega = |\mathbf{v}_h| / r$.

Si el tiro calculado chocaría con la red, una ayuda de puntería primero lo alarga manteniendo
la rapidez horizontal y, si no basta, aumenta $T$ para darle más parábola
(código~\ref{lst:tiro}).

\begin{lstlisting}[language=C++, caption={Velocidad del tiro y ayuda de puntería (\texttt{src/main.cpp}).}, label={lst:tiro}]
glm::vec3 launchVelocity(const glm::vec3& from, const glm::vec3& to, float T)
{
    return (to - from) / T - 0.5f * GRAVITY * T;
}

for (int i = 0; i < 80; i++) {
    to = origin + forward * distance;
    v = launchVelocity(from, to, T);
    if (clearsNet(from, v)) break;
    if (distance < maxDistance) { distance += 0.5f; T = distance / speed; }
    else                        { T += 0.05f; }
}
\end{lstlisting}

\subsubsection{Palmera con viento: jerarquía de transformaciones}
\label{sec:palmera}
El modelo de la palmera trae cada pieza como una malla separada: el tronco, seis hojas y
tres cocos. En lugar de dibujarlo con una sola matriz, cada malla recibe su propia
transformación, organizada como un árbol (figura~\ref{fig:jerarquia}): el tronco se mece
desde su base y las hojas y los cocos son sus hijos, así que heredan ese movimiento; además,
cada hoja aletea alrededor de su punto de unión con la punta del tronco.

\begin{figure}[H]
  \centering
  \begin{tikzpicture}[
      level distance=1.6cm, sibling distance=5.2cm,
      every node/.style={draw, rounded corners, align=center, font=\small, fill=green!6},
      edge from parent/.style={draw, -{Stealth[length=2mm]}}]
    \node {Palmera\\$M_{\text{palmera}}$}
      child { node {Tronco\\$R_{\mathbf{b}}(\theta)$ sobre la base}
        child { node {Hoja $i$ ($\times 6$)\\$R(\theta_i)$ sobre $\mathbf{p}_i$} }
        child { node {Cocos ($\times 3$)\\sin movimiento propio} } };
  \end{tikzpicture}
  \caption{Jerarquía de transformaciones de la palmera.}
  \label{fig:jerarquia}
\end{figure}

La matriz de cada hoja es la composición de la jerarquía:
\begin{equation}
  M_{\text{hoja},i} = M_{\text{palmera}}\; R_{\mathbf{b}}\!\left(\theta(t)\right)\;
  T(\mathbf{p}_i)\, R\!\left(\theta_i(t)\right)\, T(-\mathbf{p}_i),
\end{equation}
donde $R_{\mathbf{b}}$ es la rotación del tronco alrededor de su base $\mathbf{b}$ y
$\mathbf{p}_i$ el punto de unión de la hoja, que se calcula al cargar como el vértice de la
hoja más cercano a la punta del tronco. El eje de aleteo de cada hoja es horizontal y
perpendicular a su dirección. Los ángulos se obtienen con sumas de senos de frecuencias que
no son múltiplos entre sí, para que el viento no se vea periódico:
\begin{align}
  \theta(t)   &= 2.5^\circ \sin(0.9\,t) + 1^\circ \sin(2.3\,t + 1), \\
  \theta_i(t) &= 6^\circ \sin(2.5\,t + \varphi_i) + 2^\circ \sin(5.3\,t + 0.7\,\varphi_i),
\end{align}
con un desfase $\varphi_i = 1.7\,i$ distinto para cada hoja.

\subsubsection{Fuego de la fogata}
\label{sec:fuego}
La fogata tiene dos llamas (una exterior y una central). Cada una se escala desde su base
con un ruido suave $n(t) \in [-1, 1]$, formado por tres senos de frecuencias distintas:
\begin{equation}
  n(t) = 0.5\sin(7.3\,t + s) + 0.3\sin(13.1\,t + 2.1\,s) + 0.2\sin(23.7\,t + 3.7\,s),
\end{equation}
con una semilla $s$ distinta por llama. La altura se escala por $1 + 0.25\,n(t)$ y el ancho
por $1 + 0.08\,n'(t)$. Para que el fuego se vea encendido aunque no le dé el sol, se agregó al
\emph{fragment shader} un término de emisión que suma el color propio de la superficie,
\begin{equation}
  C = (k_a + k_d)\,C_{\text{base}} + C_{\text{esp}} + e(t)\, C_{\text{base}},
  \qquad e(t) = 0.45 + 0.2\,n(t),
\end{equation}
y que parpadea con el mismo ruido que la altura de la llama.

\begin{figure}[H]
  \centering
  \includegraphics[width=\linewidth]{palmera_viento}
  \caption{Palmera con viento en cuatro instantes ($t = 0$, $0.9$, $1.8$ y $2.7$ s).}
  \label{fig:palmera}
\end{figure}

\begin{figure}[H]
  \centering
  \includegraphics[width=0.7\linewidth]{fuego_parpadeo}
  \caption{Parpadeo del fuego: la llama en $t = 0$ y en $t = 0.14$ s.}
  \label{fig:fuego}
\end{figure}

\pendiente{documentar las otras dos animaciones procedurales.}

\subsection{Iluminación y texturizado}
La escena usa iluminación de Phong \cite{phong1975} con una luz direccional que representa
al sol, con componentes ambiental, difusa y especular. El fondo es un \emph{cube map}
generado a partir de la fotografía panorámica \emph{Spiaggia di Mondello} \cite{mondello},
rotada para que el sol de la foto coincida con la dirección de la luz de la escena. Los
personajes animados usan la misma iluminación que el resto de la escena.

\pendiente{segunda técnica de iluminación (por ejemplo, Fresnel para materiales traslúcidos y
metálicos), materiales opacos, plásticos, traslúcidos y metálicos, y texturizado de la arena.}

\subsection{Modelos y estadísticas}
La escena tiene 59 objetos colocados que usan 57 modelos distintos (las gradas y las bancas se
reutilizan). El inventario completo, con autor, número de piezas, vértices, caras, texturas y
origen de cada modelo, está en \texttt{docs/modelos.md}; la tabla~\ref{tab:modelos} lo resume.

\begin{table}[H]
  \centering
  \caption{Resumen de modelos registrados (al 5 de noviembre de 2026).}
  \label{tab:modelos}
  \begin{tabular}{l r r r}
    \toprule
    \textbf{Autor} & \textbf{Originales} & \textbf{Externos} & \textbf{Total} \\
    \midrule
    Parra Bello Carlos Enrique     & 33 & 0 & 33 \\
    Hernández Yukioayax Gabriel    & 32 & 8 & 40 \\
    Proyecto de teoría (Uribe, Ayala) & 4 & 0 & 4 \\
    \midrule
    \textbf{Total} & \textbf{69} & \textbf{8} & \textbf{77} \\
    \bottomrule
  \end{tabular}
\end{table}

\pendiente{actualizar al llegar a 80 modelos e incluir la tabla completa de estadísticas.}

\subsection{Interacción}
\pendiente{15 objetos interactivos.}

\subsection{Pruebas}
\label{sec:pruebas}
Para revisar cada cambio sin abrir la ventana, el programa acepta opciones de línea de
comandos que dibujan unos cuadros, guardan una captura y terminan:
\texttt{--captura} (archivo de salida), \texttt{--camara} (posición y orientación),
\texttt{--tercera} (seguir a un jugador), \texttt{--pose} (fijar una animación en un instante)
y \texttt{--simular} (ejecutar un saque o remate y adelantar la simulación, imprimiendo la
trayectoria del balón). Las figuras de este reporte se obtuvieron así.

% -----------------------------------------------------------------------------
\section{Experimentos}
% -----------------------------------------------------------------------------
\subsection{Postes dobles en la red}
Al cargar todos los modelos juntos se vieron dos pares de postes en la red: los de la cancha
del proyecto de teoría y los de la red de Gabriel (figura~\ref{fig:postes}a). Además, la red
de la cancha de teoría era solo una malla de aristas sin caras, que OpenGL no dibuja. Se
quitaron del archivo de la cancha los postes, las bandas y esa malla, y la red de la escena
quedó siendo la de Gabriel (figura~\ref{fig:postes}b).

\begin{figure}[H]
  \centering
  \begin{subfigure}{0.48\linewidth}
    \includegraphics[width=\linewidth]{exp_postes_dobles_antes}
    \caption{Antes: postes de la cancha y de la red.}
  \end{subfigure}\hfill
  \begin{subfigure}{0.48\linewidth}
    \includegraphics[width=\linewidth]{exp_postes_dobles_despues}
    \caption{Después: solo la red de Gabriel.}
  \end{subfigure}
  \caption{Experimento: postes dobles en la red.}
  \label{fig:postes}
\end{figure}

\subsection{Formato para personajes con esqueleto: FBX contra glTF}
La práctica de animación del curso carga personajes en FBX. En una práctica anterior se
comprobó que los FBX exportados por Blender 5 llegan a Assimp sin normales y el cargador del
curso falla al leerlas. Por eso se eligió glTF 2.0, que exporta normales, pesos y animaciones,
y además se agregó al cargador el cálculo de normales cuando faltan.
\textbf{Resultado:} exitoso con glTF.

\subsection{Esqueleto oculto que no se exportaba}
La primera exportación a glTF produjo un archivo con la malla pero sin esqueleto ni
animación. La causa fue que el esqueleto estaba oculto en el archivo de Blender y el
exportador omite los objetos ocultos. El script ahora hace visible el esqueleto antes de
exportar. \textbf{Resultado:} corregido.

\subsection{Pies que flotaban al caminar}
En la primera versión del ciclo de caminar, la cadera subía cuando las piernas estaban más
abiertas; como al abrir las piernas los pies suben, el personaje parecía flotar a medio paso
(figura~\ref{fig:pies}a). Se cambió el método: en cada cuadro clave se mide el punto más bajo
de los pies y se mueve la cadera para que toque el piso (figura~\ref{fig:pies}b).

\begin{figure}[H]
  \centering
  \begin{subfigure}{0.48\linewidth}
    \includegraphics[width=\linewidth]{exp_pies_flotando}
    \caption{Antes: el jugador flota a medio paso.}
  \end{subfigure}\hfill
  \begin{subfigure}{0.48\linewidth}
    \includegraphics[width=\linewidth]{exp_pies_contacto}
    \caption{Después: el pie de apoyo toca la arena.}
  \end{subfigure}
  \caption{Experimento: contacto de los pies con el piso.}
  \label{fig:pies}
\end{figure}

\subsection{Remate que pasaba por debajo de la red}
En la primera prueba del remate, el contacto con el balón quedaba a 2.0~m porque en el cuadro
del golpe el brazo ya iba hacia abajo, y el balón cruzaba el plano de la red a 1.0~m, por
debajo de ella. Se subió el brazo en ese cuadro (contacto a 2.6~m) y se agregó una ayuda de
puntería. La primera versión de la ayuda solo alargaba el tiempo de vuelo y el remate se
volvió un globo lento que subía a 3.7~m; la versión final primero alarga la distancia
manteniendo la rapidez y solo después eleva la parábola (sección~\ref{sec:tiro}). Con el
jugador a 4.5~m de la red, el remate sale a unos 13~m/s y cruza la red a 2.8~m.

% -----------------------------------------------------------------------------
\section{Resultados}
% -----------------------------------------------------------------------------
La figura~\ref{fig:general} muestra la escena completa y la figura~\ref{fig:aerea} su
distribución vista desde arriba: la cancha al centro, las gradas a los lados, la mesa de
jueces y la entrada al frente, el podio al fondo, la zona de descanso y el puesto de helados
a los costados, y el muelle, el bote y el faro hacia el mar.

\begin{figure}[H]
  \centering
  \includegraphics[width=0.9\linewidth]{vista_general}
  \caption{Vista general de la escena con la cámara libre.}
  \label{fig:general}
\end{figure}

\begin{figure}[H]
  \centering
  \includegraphics[width=0.9\linewidth]{vista_aerea}
  \caption{Distribución de la escena vista desde arriba.}
  \label{fig:aerea}
\end{figure}

\begin{figure}[H]
  \centering
  \includegraphics[width=0.9\linewidth]{camara_tercera_persona}
  \caption{Cámara en tercera persona siguiendo al jugador de México.}
  \label{fig:tercera}
\end{figure}

\begin{figure}[H]
  \centering
  \includegraphics[width=0.9\linewidth]{caminar_jugadores}
  \caption{Los cuatro jugadores en la animación de caminar.}
  \label{fig:caminar}
\end{figure}

\begin{figure}[H]
  \centering
  \includegraphics[width=\linewidth]{animaciones_keyframes}
  \caption{Cuadros de las animaciones por cuadros clave. Filas: saque visto de lado, saque de
  frente, remate de lado, y festejo de frente con la señal del árbitro.}
  \label{fig:keyframes}
\end{figure}

\begin{figure}[H]
  \centering
  \includegraphics[width=\linewidth]{tiro_parabolico}
  \caption{Tiro parabólico del balón en el saque (arriba) y el remate (abajo), en distintos
  instantes de la animación.}
  \label{fig:tiro}
\end{figure}

\begin{table}[H]
  \centering
  \caption{Avance de los requisitos del proyecto.}
  \label{tab:avance}
  \begin{tabularx}{\linewidth}{X l}
    \toprule
    \textbf{Requisito} & \textbf{Estado} \\
    \midrule
    80 modelos (al menos 80\% originales) & 77 de 80 \\
    2 modelos de cámara & Cumplido \\
    15 objetos interactivos & Pendiente \\
    Materiales opacos, plásticos, traslúcidos y metálicos & Pendiente \\
    2 técnicas de iluminación & 1 de 2 (Phong) \\
    Mapeo de ambiente & Cumplido (\emph{cube map}) \\
    10 animaciones simples & Pendiente \\
    5 animaciones por cuadros clave & Cumplido \\
    5 animaciones procedurales & 3 de 5 \\
    \bottomrule
  \end{tabularx}
\end{table}

\pendiente{discusión de resultados comparada con otros desarrollos; ventajas y desventajas de
la metodología.}

% -----------------------------------------------------------------------------
\section{Conclusiones}
% -----------------------------------------------------------------------------
\subsection*{Hernández Yukioayax Gabriel}
\pendiente{conclusión individual.}

\subsection*{Parra Bello Carlos Enrique}
\pendiente{conclusión individual.}

% -----------------------------------------------------------------------------
\section*{Uso de inteligencia artificial}
\addcontentsline{toc}{section}{Uso de inteligencia artificial}
% -----------------------------------------------------------------------------
Se usó el asistente Claude (Anthropic) como apoyo en el desarrollo: en los scripts de Python
para Blender (exportación de modelos, esqueleto de los personajes y creación de las
animaciones), en el código de carga de modelos con esqueleto, en la simulación del balón y en
la redacción de este documento. El uso de IA en los modelos se detalla en la sección
``Uso de IA'' de \texttt{docs/modelos.md}. Todo el contenido fue revisado y ajustado por los
autores.

% -----------------------------------------------------------------------------
\begin{thebibliography}{9}
\addcontentsline{toc}{section}{Referencias}
% Formato APA 7. Ordenadas alfabéticamente por autor.

\bibitem{assimp}
Assimp Team. (2024). \emph{Open Asset Import Library (Assimp)} [Software].
\url{https://github.com/assimp/assimp}

\bibitem{blender}
Blender Foundation. (2026). \emph{Blender} (Versión 5.2) [Software].
\url{https://www.blender.org}

\bibitem{devries2020}
de Vries, J. (2020). \emph{Learn OpenGL: Learn modern OpenGL graphics programming in a
step-by-step fashion}. Kendall \& Welles. \url{https://learnopengl.com}

\bibitem{halliday2014}
Halliday, D., Resnick, R., \& Walker, J. (2014). \emph{Fundamentals of physics} (10.ª ed.).
Wiley.

\bibitem{gltf2021}
Khronos Group. (2021). \emph{glTF 2.0 specification}.
\url{https://registry.khronos.org/glTF/specs/2.0/glTF-2.0.html}

\bibitem{lasseter1987}
Lasseter, J. (1987). Principles of traditional animation applied to 3D computer animation.
\emph{ACM SIGGRAPH Computer Graphics, 21}(4), 35--44.
\url{https://doi.org/10.1145/37402.37407}

\bibitem{magnenat1988}
Magnenat-Thalmann, N., Laperrière, R., \& Thalmann, D. (1988). Joint-dependent local
deformations for hand animation and object grasping. En \emph{Proceedings of Graphics
Interface '88} (pp. 26--33). Canadian Information Processing Society.

\bibitem{mondello}
Mischok, A. (s.f.). \emph{Spiaggia di Mondello} [Imagen HDRI, licencia CC0]. Poly Haven.
\url{https://polyhaven.com/a/spiaggia_di_mondello}

\bibitem{phong1975}
Phong, B. T. (1975). Illumination for computer generated pictures. \emph{Communications of
the ACM, 18}(6), 311--317. \url{https://doi.org/10.1145/360825.360839}

\end{thebibliography}

\end{document}
